\chapter{ABAQUS input files}
\label{ch:input}

The model is defined in an ABAQUS/CAE input file (\code{*.inp}). The readers in
\code{subroutines/ReadInputFileAbaqus} are not a general ABAQUS parser. They
expect the structure that ABAQUS/CAE writes, including several of its comment
lines. The safest approach is to build the model in ABAQUS/CAE (one part, one
instance, 2D planar, CPS4/CPE4 or CPS8/CPE8 elements) and to start from one of
the examples in \code{code/InputFilesAbaqus}.

\section{Reading sequence}

\begin{enumerate}
  \item \code{Adjust_Set_names_input_file_ABAQUS(filename)} writes
    \code{<name>_adjusted.inp} to the current folder. In it, every \code{Set-}
    after \code{End Instance} is renamed to \code{PickedSet-}, and the node/element
    set pairs named \code{Pressure} and \code{Traction} are numbered
    (\code{Pressure1}, \code{Pressure2}, \dots) in order of appearance.
  \item \code{Read_mesh_ABAQUS}: nodes $(x,y)$ and element connectivity.
  \item \code{Read_material_ABAQUS}: materials and their element sets.
  \item \code{Read_step_ABAQUS}: step names, types and the time period of
    \code{*Dynamic} steps.
  \item \code{Read_load_ABAQUS}: concentrated forces, edge loads, body forces and
    gravity.
  \item \code{Read_boundary_condition_ABAQUS}: Dirichlet conditions and nodes of
    interest.
  \item The adjusted file is deleted.
  \item \code{pSpectralElement_NodesTopology} and \code{pSpectralElement_BCs} lift
    the mesh and the edge-based sets to degree $p$.
\end{enumerate}

\section{Requirements and conventions}

\paragraph{Mesh.}
\begin{itemize}
  \item Only the first \code{*Node} block is read, and only $x$ and $y$. Node
    labels must be numbered consecutively from 1 (checked by the reader).
  \item Element types ending in \code{4}/\code{4R} give 4-node elements, types
    containing \code{8} give 8-node elements. One element type per model.
  \item For 8-node meshes the corner nodes must be numbered before the mid-side
    nodes (\code{pSpectralElement_NodesTopology} keeps nodes \code{1..nCorner} and
    stops with an error otherwise).
  \item The element block must be followed directly by \code{*Nset},
    \code{*End Part} or \code{*End Instance}.
\end{itemize}

\paragraph{Sections and materials.}
\begin{itemize}
  \item \code{*Solid Section, elset=<set>, material=<mat>}, preceded by a
    \code{** Section:} comment.
  \item A \code{** MATERIALS} header; the material block ends with the
    \code{** ----} line that CAE writes.
  \item \code{*Density}, \code{*Elastic} (isotropic \code{E, nu},
    \code{ENGINEERING CONSTANTS} or \code{ORTHO}) and \code{*Damping}.
    \code{*Damping} only switches Rayleigh damping on for the element set; the
    parameters come from \code{omega_vec} and \code{zeta_vec}.
\end{itemize}

\paragraph{Sets.}
\begin{itemize}
  \item Sets used by boundary conditions and loads must be defined in the
    \emph{assembly} (after \code{*End Instance}), as CAE does. Otherwise the
    boundary condition reader stops with the error ``Node set \dots\ was not
    found''. Set names are matched exactly (\code{Set-1} is not \code{Set-10}).
  \item Node and element sets with \code{generate} are expanded by the readers
    (also for \code{*Cload}).
  \item Boundary-condition set names must contain \code{Set} (\code{Set-1},
    \code{_PickedSet3}, \dots).
  \item A set on an \emph{edge} needs a node set \emph{and} an element set with the
    same name (CAE creates both when edges are picked). Only then can
    \code{pSpectralElement_BCs} add the high-order nodes of that edge. Point
    constraints only need a node set.
  \item Edge loads: define node set and element set as \code{Pressure} (pressure
    loads) or \code{Traction} (surface tractions), in the order of the loads in
    the step.
\end{itemize}

\paragraph{Step, boundary conditions and loads} (only data after
\code{** STEP:} is read):

\begin{xltabular}{\linewidth}{@{}P{5.6cm}P{2.6cm}L@{}}
  \toprule
  ABAQUS & internal \code{type} & notes \\
  \midrule
  \endhead
  \code{*Boundary} with \code{Set-1, ENCASTRE} or \code{PINNED} & \code{fixed} & $u_x = u_y = 0$ \\
  \code{Set-1, XSYMM} / \code{YSYMM} & \code{XSYMM} / \code{YSYMM} & \\
  \code{Set-1, 1, 1} & \code{Ux 0} & homogeneous \\
  \code{Set-1, 1, 1, 0.001} & \code{Ux} & prescribed value 0.001, applied by elimination \\
  \code{Set-1, 1, 1} and \code{Set-1, 2, 2} (two lines, as CAE writes U1 + U2) & \code{UxUy 0} & \\
  \code{*Boundary, amplitude=...} & \code{Ux dt}, \code{Uy dt} & time-dependent; time function from \code{SignalType} \\
  \code{*Boundary, user} & \code{Ux fx}, \code{Uy fx} & spatially varying, \code{SpatialDisplacementFunction.m} \\
  \code{Set-5, 6, 6} & \code{noi} & node of interest for the output of the dynamic analysis; not a constraint \\
  \code{*Cload}: \code{Set, 1, F} & \code{f_x} / \code{f_y} & full value at every node of the set; \code{amplitude=} gives \code{f_x_dt} \\
  \code{*Dsload}: \code{Pressure, P, p} & \code{Pressure} & along the inward normal \\
  \code{*Dsload}: \code{Pressure, PNU, a} & \code{Pressure fx} & $p(x,y)$ from \code{SpatialLoadingFunction.m} \\
  \code{*Dsload}: \code{Traction, TRVEC, a, dx, dy, dz} & \code{Traction} & traction $a\,(d_x, d_y)$ \\
  \code{*Dsload}: \code{Traction, TRVECNU, a, dx, dy, dz} & \code{Traction fx} & $a\,(d_x, d_y)\,f(x,y)$ from \code{SpatialLoadingFunction.m} \\
  \code{*Dload} with \code{BX}/\code{BY} (CAE comment \code{Body force}) & \code{Body_force} & \\
  \code{*Dload}: \code{..., GRAV, g, dx, dy} & gravity & all elements \\
  \bottomrule
\end{xltabular}

\code{*Amplitude} tables are ignored: the time history of every time-dependent
load is the signal selected by \code{SignalType} in \code{Start_Sim.m}.
\code{SpatialLoadingFunction.m} contains the loads of the Lee \& Bathe cubic
patch test: \code{Pressure fx} gives $a\cdot 120\,y$, \code{Traction fx} gives
$a\,(6 - 6y^2)$.

\section{Data structures returned by the readers}

\begin{xltabular}{\linewidth}{@{}P{4.6cm}L@{}}
  \toprule
  variable & layout \\
  \midrule
  \code{nodes} & \code{[nNodes x 2]} coordinates \\
  \code{topology} & \code{[nElem x 4]} or \code{[nElem x 8]} connectivity \\
  \code{material_properties{i,1:5}} & density, dielectric (legacy), elastic
    constants (\code{[E nu]} or 9 orthotropic stiffness values), piezoelectric
    (legacy), damping \\
  \code{material_elset{i}} & elements of material \code{i} \\
  \code{cload(k)}, \code{dsload(k)} & struct arrays with the fields
    \code{set_name}, \code{amplitude}, \code{type}, \code{nset}, \code{elset},
    \code{DirectionVector}, \code{step_name} (\code{cload} has no \code{elset}
    and no \code{DirectionVector}) \\
  \code{body_force(k)}, \code{gravity(k)} & \code{amplitude = [bx; by]} or
    \code{[g; dx; dy]}, \code{elset} \\
  \code{boundary_conditions(k)} & column struct array with the fields
    \code{set_name}, \code{amplitude}, \code{type}, \code{nset}, \code{elset},
    \code{step_name}, in reverse order of appearance \\
  \bottomrule
\end{xltabular}

\section{Example: Lee \& Bathe quadratic patch test}
\label{sec:leebathe}

The file \code{PatchTestLeeQuadratic_angular.inp} contains a plate
$[0,10]\times[-1,1]$ with two 4-node elements and angular distortion,
$u_x = 0$ on the left edge (\code{Set-1}, node set and element set), $u_y = 0$
at $(0,-1)$ (\code{Set-2}) and the pressure $p = 120\,y$ on the right edge
(\code{Pressure}, \code{PNU}); see \cref{fig:leebathe}. The exact solution is
pure bending (plane stress, $\sigma_x = a y$):
\begin{equation}
  u_x = \frac{a\,x\,y}{E},\qquad
  u_y = -\frac{a}{2E}\left(x^2 + \nu y^2\right) + \frac{a\,\nu}{2E}.
\end{equation}
For $p\ge2$ the unsymmetric spectral element reproduces it to machine precision
on all three meshes (\code{tests/integration/LeeBathePatchTest.m},
\code{tests/integration/AbaqusInputTest.m}).

\begin{figure}[htbp]
  \centering
  % Lee & Bathe (1993) quadratic patch test: plate [0,10] x [-1,1],
% boundary conditions, load and the three meshes of code/InputFilesAbaqus
\begin{tikzpicture}[x=0.62cm, y=0.62cm, every node/.style={font=\small}]
  % ---- model with boundary conditions and load ----
  \begin{scope}
    \draw[thick, fill=black!4] (0,-1) rectangle (10,1);
    \node[above=2pt] at (5,1) {$E=70\,\mathrm{GPa}$, $\nu=0.3$, plane stress};
    % u_x = 0 on the left edge: rollers on a vertical wall
    \foreach \y in {-0.5,0,0.5,1} {
      \draw (0,\y) -- (-0.35,\y+0.17) -- (-0.35,\y-0.17) -- cycle;
    }
    \draw[thick] (-0.48,-0.75) -- (-0.48,1.25);
    \node[left] at (-0.6,0.25) {$u_x=0$};
    % u_y = 0 (and u_x = 0) at (0,-1): pin
    \draw (0,-1) -- (-0.22,-1.42) -- (0.22,-1.42) -- cycle;
    \draw[thick] (-0.35,-1.42) -- (0.35,-1.42);
    \node[below] at (0,-1.45) {$u_y=0$};
    % linearly varying pressure p = 120 y on the right edge: compression
    % (pointing inwards) for y > 0, tension (pointing outwards) for y < 0
    \foreach \y in {0.25,0.5,0.75,1} {
      \draw[-{Stealth[length=1.4mm]}] (10+1.1*\y,\y) -- (10,\y);
      \draw[-{Stealth[length=1.4mm]}] (10,-\y) -- (10+1.1*\y,-\y);
    }
    \draw (11.1,1) -- (10,0) -- (11.1,-1);
    \node[right, align=left] at (11.4,0)
      {$p=120\,y$\\(\code{PNU}, \code{SpatialLoadingFunction})};
  \end{scope}
  % ---- meshes (node coordinates of the input files) ----
  \begin{scope}[shift={(0,-4.2)}]
    \draw[thick] (0,-1) rectangle (10,1);
    \draw[thick] (5,-1) -- (5,1);
    \node[left] at (-0.3,0) {undistorted};
  \end{scope}
  \begin{scope}[shift={(0,-7.2)}]
    \draw[thick] (0,-1) rectangle (10,1);
    \draw[thick] (4,-1) -- (6,1);
    \node[left] at (-0.3,0) {angular};
  \end{scope}
  \begin{scope}[shift={(0,-10.2)}]
    \draw[thick] (0,-1) rectangle (10,1);
    % parabolic edge through (5,1), (3.5,0), (5,-1): quadratic Bezier with
    % control point (2,0), written as the equivalent cubic Bezier
    \draw[thick] (5,1) .. controls (3,1/3) and (3,-1/3) .. (5,-1);
    \foreach \p in {(5,1),(3.5,0),(5,-1)} \node[node dot] at \p {};
    \node[left] at (-0.3,0) {mid-edge};
  \end{scope}
\end{tikzpicture}

  \caption{Lee \& Bathe quadratic patch test: model, boundary conditions and the
    undistorted, angular and mid-edge (curved) meshes of
    \texttt{code/InputFilesAbaqus}.}
  \label{fig:leebathe}
\end{figure}

The input files were reconstructed from the result workspaces that came with the
code (\code{tests/fixtures/make_abaqus_inputs.m}); they reproduce the stored mesh,
loads and boundary conditions exactly.
